\section{Threat Model and Attack Taxonomy}\label{sec:threat}
% Written in P8 (framework-independent). Sources: configs/eda.yaml (rule, fixed before any metric was computed),
% configs/attack_groups.yaml, results/tables/attack_signatures.csv, periodicity_summary.csv, capture_lineage.csv;
% decisions K-007, K-009 ... K-014. No numeric literal in this file: every number is a macro (generated/numbers.tex).

\subsection{Adversary and defender}
A CAN frame carries an identifier, which also sets its arbitration priority, and up to \NPayloadBytes{} payload
bytes. The protocol does not authenticate the sender, so any node that can transmit on the bus can emit a frame
with an arbitrary identifier and payload.
% TODO-CITE: CAN standard (ISO 11898) and the absence of sender authentication; in-vehicle attack surface
% (compromised ECU, diagnostic port).
We consider an adversary with this capability and restrict the attack classes to those that the two datasets
contain: \emph{flooding} with a high-priority identifier (identifier \DoSTarget{} in Car-Hacking), \emph{fuzzing}
with random identifiers and payloads (Car-Hacking) or an all-ones payload on many identifiers (ROAD),
\emph{fabrication} (spoofing), where frames with a valid identifier are injected next to those of the legitimate
sender, and \emph{masquerade}, where the injected frames replace those of the legitimate sender, so that the
identifier keeps its normal schedule and only its payload changes.
% TODO-CITE: masquerade / suspension attacks and how a legitimate ECU can be silenced (e.g. bus-off).

The defender is an intrusion detection system on an edge ECU that passively reads every frame and labels each window
of \WinSize{} consecutive frames as attack or normal. The model is trained offline, compressed by post-training
quantization and deployed with a decision threshold that is fixed before deployment from attack-free validation
traffic (Section~\ref{sec:method}). Prevention, attacks below the frame level and adaptive adversaries that know
the detector are out of scope.

\subsection{A mechanism-based taxonomy}
The intuition behind this study is that CAN attacks differ in how visible they are in the traffic structure. We
define the two groups by mechanism rather than by attack name:
\begin{itemize}
\item \emph{coarse} attacks change the timing or the identifier distribution of the bus, which includes every
attack that adds frames;
\item \emph{stealthy} attacks preserve both, so that their signature lies only in the payload.
\end{itemize}
The name-based grouping this study started from (flooding and fuzzing coarse; spoofing, malfunction and masquerade
stealthy) is an assumption about traffic, not a measurement. We therefore assign each attack type with four traffic
criteria, computed between the attack period and the normal period of the same captures. A type is coarse if any
criterion fires and stealthy otherwise:
\begin{itemize}
\item[C1] the share of attack-period frames whose identifier never occurs in attack-free reference traffic
exceeds \COneThr;
\item[C2] the median Jensen--Shannon divergence between the identifier frequencies of \CTwoFrames-frame attack
windows and those of the normal period exceeds the \CTwoPct{}th percentile of a null distribution drawn
\CTwoDraws{} times from attack-free traffic (equal window sizes keep the sample-size dependence of the divergence
identical);
\item[C3] the frame rate of the target identifier changes by more than a factor of \CBand{} in either direction;
\item[C4] the median inter-arrival time of the target identifier changes by more than a factor of \CBand{}.
\end{itemize}
Attacks without a single target identifier (fuzzing) are judged on C1 and C2 only. In Car-Hacking, injected frames
less than \BurstGapS{}~s apart form a burst and the bursts form the attack period; in ROAD the attack period is the
injection interval of the capture metadata. Bus silences longer than \RecGapS{}~s are recording gaps and are
excluded from rates and inter-arrival times. The rule and its thresholds were fixed before any of these statistics
was computed, and it uses traffic statistics only: no model output enters the grouping, which keeps the later
comparison of quantization losses between the groups from being circular.

\subsection{Empirical grouping}
Table~\ref{tab:signatures} lists the criteria for all \NTypes{} attack types (Fig.~\ref{fig:signatures} in the
appendix shows them graphically). CAN traffic is strongly periodic, which is what makes timing informative: in
Car-Hacking all \PeriodicCH{} of \IdsCH{} identifiers have an inter-arrival-time coefficient of variation below
\CVThr{} (median \MedianCVCH), and in ROAD \RobustPeriodicROAD{} of \IdsROAD{} identifiers meet the same bound with
a robust spread measure.

Flooding and fuzzing change the identifier distribution: in Car-Hacking \DoSUnseenPct{} (DoS) and
\FuzzyUnseenPct{} (fuzzy) of the attack-period frames carry an unseen identifier, and the window divergence
(\DoSJSD{} and \FuzzyJSD) is far above the null percentile (\NullQCH). ROAD fuzzing fires only C1
(\RoadFuzzUnseenPct): its divergence (\RoadFuzzJSD) stays below the ROAD null (\NullQROAD).

Spoofing and fabrication use valid identifiers and add no unseen identifier, but they are not stealthy. The
Car-Hacking spoofing attacks multiply the target-identifier rate by \SpoofRate{} and scale its median inter-arrival
time by \SpoofIAT{}, enough to move the identifier distribution as well (divergence \SpoofJSD). The ROAD fabrication
attacks double the target rate (ratio \FabRate) and shorten its median inter-arrival time (ratio \FabIAT), while
the identifier distribution changes less than the noise floor (divergence \FabJSD{} against \NullQROAD): a detector
that only looked at identifier frequencies would miss them, but the timing of the target identifier gives them
away.

Masquerade attacks leave every traffic statistic at its normal value: the target rate ratio is \MasqRate, the
inter-arrival-time ratio \MasqIAT, the divergence at most \MasqJSDMax{} and no unseen identifier occurs. Their
signature is in the payload only; the largest shift of a payload-byte median between attack and normal period is
\MasqByteShift{} byte units, depending on the attack.

The rule yields \NCoarse{} coarse and \NStealthy{} stealthy types (\NCoarseCaptures{} and \NStealthyCaptures{}
captures). It contradicts the name-based grouping in \NPlanContradict{} types: every spoofing and fabrication
attack is coarse. All stealthy types are ROAD masquerade attacks; Car-Hacking contains none.

\begin{table}[t]
\centering
\footnotesize
\setlength{\tabcolsep}{3pt}
\caption{Traffic criteria that define the attack groups (exploratory analysis, before any model was trained). C1: share of attack-period frames with an ID unseen in attack-free traffic (coarse if $>0.01$). C2: median Jensen--Shannon divergence of the ID distribution of 1{,}000-frame attack windows against the normal period of the same capture (coarse if above the 99th percentile of an attack-free null: CH 0.00059, ROAD 0.077). C3, C4: target-ID rate and median inter-arrival-time ratio, attack/normal period (coarse if outside $[1/1.5,\,1.5]$; not defined without a single target ID). Bold: criterion fires. A type is coarse if any criterion fires.}
\label{tab:signatures}
\begin{tabular}{@{}lrrrrl@{}}
\toprule
Attack type & C1 & C2 & C3 & C4 & Group \\
\midrule
\multicolumn{6}{@{}l}{\textit{Car-Hacking}} \\
DoS & \textbf{0.545} & \textbf{0.349} & \textbf{$\infty$} & \textbf{$\infty$} & coarse \\
fuzzy & \textbf{0.379} & \textbf{0.226} & -- & -- & coarse \\
gear spoofing & 0.000 & \textbf{0.111} & \textbf{6.47} & \textbf{0.11} & coarse \\
RPM spoofing & 0.000 & \textbf{0.111} & \textbf{7.06} & \textbf{0.11} & coarse \\
\midrule
\multicolumn{6}{@{}l}{\textit{ROAD}} \\
fuzzing & \textbf{0.066} & 0.040 & -- & -- & coarse \\
correlated signal & 0.000 & 0.010 & \textbf{2.00} & \textbf{0.41} & coarse \\
max speedometer & 0.000 & 0.011 & \textbf{2.00} & \textbf{0.36} & coarse \\
reverse light off & 0.000 & 0.010 & \textbf{2.00} & \textbf{0.40} & coarse \\
reverse light on & 0.000 & 0.010 & \textbf{2.00} & \textbf{0.39} & coarse \\
correlated signal (masq.) & 0.000 & 0.006 & 1.00 & 1.02 & stealthy \\
max speedometer (masq.) & 0.000 & 0.006 & 1.00 & 1.00 & stealthy \\
reverse light off (masq.) & 0.000 & 0.006 & 1.00 & 1.00 & stealthy \\
reverse light on (masq.) & 0.000 & 0.005 & 1.00 & 1.00 & stealthy \\
max coolant temp & 0.000 & 0.006 & \textbf{2.03} & \textbf{0.50} & coarse \\
max coolant temp (masq.) & 0.000 & 0.006 & 1.03 & 1.14 & stealthy \\
\bottomrule
\end{tabular}
\end{table}


\subsection{Paired fabrication--masquerade design}
Each of the \NMasqCaptures{} ROAD masquerade captures is derived from a fabrication capture: it is the same log
with the conflicting legitimate target-identifier frames removed from the injection interval. We verified this
lineage on the data: \MasqInSourcePct{} of the frames of every masquerade capture occur in its source capture
(same relative time, identifier and payload), every removed frame is a legitimate target-identifier frame inside
the injection interval, and both versions keep the same injected frames. A fabrication capture and its masquerade
version therefore share the vehicle, the driving state and the injected payload values; they differ only in the
timing trace that the extra frames leave (the doubled target rate). Comparing the quantization loss within such a
pair isolates the contribution of that timing cue, with the confounders of a between-type comparison held fixed.
This pairing is the primary design for the comparison of the two groups; the comparison of the groups as a whole
is secondary (Section~\ref{sec:setup}). Because the two members of a pair contain the same injected frames, they
are always assigned to the same split (Section~\ref{sec:method}).

The masquerade captures are synthetic in one respect: in a real masquerade attack the legitimate ECU has to be
silenced, which may leave timing side effects that removing frames from a log does not reproduce. The stealthy
group thus rests on one vehicle, one dataset and post-processed captures (Section~\ref{sec:limitations}).

\subsection{Why the taxonomy matters for quantization}
A detector can recognise a coarse attack from gross structure, such as unseen identifiers or a doubled frame rate,
which reduced numerical precision is unlikely to erase. A stealthy attack is visible only as a payload deviation,
which for some attacks amounts to a few byte units; this is the kind of fine difference that coarser arithmetic
could suppress. This motivates the question whether post-training quantization costs more detection performance on
stealthy than on coarse attacks, and whether an aggregate score would hide such a cost, since coarse attacks
dominate both datasets (Table~\ref{tab:dataset}). To make sure that a measured loss comes from the quantized model
and not from the way the input is scaled, the input encoding keeps identifier codes and payload bytes exactly
representable on the INT8 input grid (Section~\ref{sec:method}).
